%% =========================================================================
%% MorphoSafety -- working draft in GECCO format (outline + bullet points)
%% Started 2026-10-12. Serves the end-of-October report first, and a later
%% GECCO 2027 submission. Every bullet is content to expand into prose.
%%
%% GECCO format (2026 instructions; the 2027 call is not out yet):
%%   * ACM sigconf template (acmart >= 2.16, Libertine fonts), US Letter
%%   * at most 8 pages excluding references; abstract at most 200 words
%%   * double-blind: no names, affiliations or acknowledgments, and own
%%     work cited in the third person -> switch to the submission line below
%%   * GECCO 2027: Krakow, Poland, 12-16 July 2027 (third-party listing);
%%     deadline expected late January 2027 (2026: 19 Jan abstract, 26 Jan paper)
%%   * Track: Complex Systems (CS), which covers evolutionary robotics
%%     (evolved controllers, morphologies and sensors)
%%
%% Status tags used in the text:
%%   [HAVE]  result in hand (preliminary unless stated)
%%   [RUN]   experiment still to run
%%   [CHECK] number or claim to verify / re-measure
%%   [WRITE] prose, figure or diagram still to produce
%%   [OPTIONAL] stretch goal
%%
%% Build: tectonic main.tex   |   latexmk -pdf main.tex   |   upload to Overleaf
%% =========================================================================

\documentclass[sigconf,review]{acmart}
%% Submission version (anonymous, line numbers):
%% \documentclass[sigconf,review,anonymous]{acmart}

%% Draft-only front matter: no copyright block or ACM reference box.
\settopmatter{printacmref=false, printfolios=true}
\setcopyright{none}
\renewcommand\footnotetextcopyrightpermission[1]{}
\acmConference[GECCO '27]{Genetic and Evolutionary Computation Conference}{July 12--16, 2027}{Krak\'{o}w, Poland}
\acmYear{2027}
\copyrightyear{2027}
\acmDOI{}
\acmISBN{}

%% ---- draft helpers -------------------------------------------------------
\newcommand{\statustag}[2]{\textcolor{#1}{\textsf{\scriptsize\bfseries[#2]}}}
\newcommand{\HAVE}{\statustag{teal!70!black}{HAVE}\hspace{0.3em}}
\newcommand{\RUN}{\statustag{red!75!black}{RUN}\hspace{0.3em}}
\newcommand{\CHECK}{\statustag{orange!85!black}{CHECK}\hspace{0.3em}}
\newcommand{\WRITE}{\statustag{blue!70!black}{WRITE}\hspace{0.3em}}
\newcommand{\OPT}{\statustag{gray}{OPTIONAL}\hspace{0.3em}}
\newcommand{\sectarget}[1]{\par\noindent{\color{gray}\sffamily\scriptsize Target length: #1}\par\smallskip}
\newcommand{\figbox}[2]{\fbox{\parbox[c][#1][c]{0.95\linewidth}{\centering\sffamily\footnotesize #2}}}

\begin{document}

%% Alternative titles:
%%   "Safety Is Also a Body Problem: Constrained Co-Design of Morphology and Control"
%%   "Evolving Safer Bodies: Constraint-Aware Co-Design of Legged Robots"
\title[Morphology as a Lever for Safe RL]{Morphology as a Lever for Safe Reinforcement Learning: Constrained Co-Design of Body and Controller}

\author{Author One}
\affiliation{%
  \institution{Institution (TODO)}
  \city{City}
  \country{Country}}
\email{author@example.org}

\author{Author Two}
\affiliation{%
  \institution{Institution (TODO)}
  \city{City}
  \country{Country}}
\email{author2@example.org}

\renewcommand{\shortauthors}{Anonymous et al.}

\begin{abstract}
\noindent\textit{Final version: one paragraph, at most 200 words. Content, in order:}
\begin{itemize}
\item \textbf{Problem.} Safe reinforcement learning (RL) optimises the controller for a fixed body, but whether a fast and safe behaviour exists, and whether it can be learned, depends on the body.
\item \textbf{Approach.} Constrained co-design: an evolved Gaussian-mixture search distribution over ant limb and torso dimensions, a morphology-conditioned PPO-Lagrangian controller trained across a population of bodies simulated side by side on one GPU, one Lagrange multiplier per mixture component, and feasibility-first ranking of designs.
\item \textbf{Result 1.} On a hazard-field crossing task the co-designed ant is about 55\% faster than the nominal ant at a comparable safety cost (1.89 vs 1.20--1.22\,m/s). \CHECK{} final numbers.
\item \textbf{Result 2.} Long legs and a small torso evolve consistently; shin length tracks the hazard statistics. Ranking designs by speed alone selects bodies 5$\times$ over the budget when the constraint is hard to meet, so the outer loop needs constraint handling too.
\item \textbf{Result 3.} An open GPU stack for safe RL and safe co-design on MuJoCo Warp: 118--133k steps/s on one RTX 3090, and a 2-billion-step co-design run in about 4\,h.
\item \textbf{Takeaway.} Morphology is a lever for safety, not only for performance.
\end{itemize}
\end{abstract}

%% TODO generate the CCSXML block at https://dl.acm.org/ccs and paste it here.
\ccsdesc[500]{Computing methodologies~Reinforcement learning}
\ccsdesc[500]{Computer systems organization~Evolutionary robotics}
\ccsdesc[300]{Computing methodologies~Physical simulation}

\keywords{morphology-control co-design, safe reinforcement learning, constrained Markov decision processes, constraint handling, evolutionary robotics, legged locomotion, GPU simulation}

\maketitle

%% =========================================================================
\section*{Draft status and plan (delete before submission)}
\begin{itemize}
\item \textbf{What this is.} An outline for the end-of-October report and a later GECCO 2027 submission (Complex Systems track). Tags: \HAVE result in hand, \RUN experiment to run, \CHECK number or claim to verify, \WRITE prose or figure to produce, \OPT stretch.
\item \textbf{Claim we can make today.} The co-designed ant is about 55\% faster than the nominal ant at a comparable episode cost (one seed). \textbf{Claim we want.} Co-design moves the safety--performance frontier (faster at equal cost \emph{and} cheaper at equal speed), and the safety constraint changes which body evolves.
\item \textbf{Every \HAVE number is preliminary.} They come from the pre-Warp stack (git tag \texttt{pre-warp}: MJX-JAX, MuJoCo 3.3.2, solver capped at 4 iterations / 8 line-search steps) with one seed. The same policy runs 4.5\% slower on the current stack, so regenerate every table on MuJoCo Warp with the full solver and never mix the two stacks in one table.
\item \textbf{Experiment queue.} Costs in RTX 3090 GPU-hours on Warp: about 1\,h per 300M-step single-body run and 5\,h per 2B-step co-design run.
  \begin{enumerate}
  \item \emph{P1, needed for any claim (report):} nominal ant with PPO-Lagrangian at $d=2$, 3 seeds (3\,h); co-design at $d=2$ on the lattice, 3 seeds (15\,h); the evolved body retrained from scratch at $d=2$, 3 seeds (3\,h); unconstrained co-design, 3 seeds (15\,h), then its body trained with the constraint (3\,h).
  \item \emph{P2, method contribution:} implement feasibility-first ranking; on random fields, speed-only vs feasibility-first, 3 seeds each (30\,h); shared vs per-component $\lambda$ (15\,h).
  \item \emph{P3, headline figure:} budget sweep $d \in \{0, 5, 25\}$ for the nominal ant and for co-design (20--45\,h).
  \item \emph{Optional:} the lasers + mines task (body pressures that oppose each other); same-task throughput against the original mjx-safety-gym, Safety-Gymnasium (CPU) and CRAX.
  \end{enumerate}
  About 100--125 GPU-hours in all; roughly a week of wall clock with 4--6 cluster jobs in parallel.
\item \textbf{Open decisions.}
  \begin{enumerate}
  \item Keep actuator rescaling (search shape at constant torque margin), or fix the gear or add an energy cost so that size is not free.
  \item Widen the design box: 4 of 7 genes sit at the 1.4 edge, and Schaff et al. used 0.5--1.5.
  \item Envs per GPU: stay at 1024 (as all current runs) or move to 4096. 4096 is 1.75$\times$ faster, but co-design iterations must lengthen and learning at that size is unverified.
  \item Scope of the simulator contribution (Sec.~\ref{sec:res-sim}).
  \end{enumerate}
\end{itemize}

%% =========================================================================
\section{Introduction}
\label{sec:intro}
\sectarget{1 page including Fig.~\ref{fig:teaser}}
\begin{itemize}
\item \textbf{Hook.} Legged robots increasingly have to move quickly through places where some contacts are unacceptable: crop rows they must not trample, clutter they must not disturb, hazards they must not step on. Safe RL formalises this as a constrained Markov decision process (CMDP): maximise return subject to a bound on expected cost~\cite{altman1999constrained,achiam2017cpo,ray2019benchmarking}.
\item \textbf{The overlooked variable.} Safe RL optimises the \emph{controller} for a \emph{fixed} body. Yet the safety--performance trade-off a robot can reach belongs to body and controller together: morphology shapes which behaviours are possible, robust and easy to learn~\cite{pfeifer2006body,bongard2011morphological}. Example from our task: legs long enough to clear a hazard turn ``weave around it'' into ``step over it''.
\item \textbf{Research question.} Can optimising morphology improve the safety--performance trade-off beyond what constrained policy optimisation reaches on the nominal body? That is, is morphology a lever for safety and not only for performance?
\item \textbf{Gap.}
  \begin{itemize}
  \item Co-design work optimises reward only~\cite{sims1994evolving,ha2019reinforcement,schaff2019jointly,wang2019nge,yuan2022transform2act,gupta2021embodied}.
  \item Where safety appears in co-design, it is a penalty inside the reward~\cite{codesign2025notfree}, the task itself (fall recovery~\cite{yue2025humanoid}), or a design specification such as velocity, power or payload~\cite{glamdring2026}.
  \item Safe-RL work keeps the body fixed~\cite{ray2019benchmarking,ji2023safetygymnasium,kim2024constraints,chanesane2024cat}.
  \item A systematic citation search (Sec.~\ref{sec:related}) found no work that trains the controller with constrained RL inside a design loop. \CHECK{} rerun before submission.
  \end{itemize}
\item \textbf{Why it is not just ``add a cost term''.} There are two coupled constraint problems:
  \begin{enumerate}
  \item \emph{Inner.} The budget must hold for every body, not on average over the population the shared policy trains on; otherwise a reckless body can hide behind safe ones.
  \item \emph{Outer.} Ranking designs by speed assumes the inner loop has already met the budget. When it has not, the search selects fast, unsafe bodies, which we observe (Sec.~\ref{sec:res-outer}). This is the classic constraint-handling problem of evolutionary computation~\cite{deb2000efficient,runarsson2000stochastic,coello2002theoretical}, with one twist: feasibility is a noisy estimate that depends on a controller that is still learning.
  \end{enumerate}
\item \textbf{Compute.} Population-based co-design needs billions of simulation steps. GPU simulation makes that practical~\cite{freeman2021brax,makoviychuk2021isaac,zakka2025playground}; we give every body its own lanes in one GPU batch.
\item \textbf{Contributions.}
  \begin{enumerate}
  \item \emph{A safe co-design method.} A constrained bilevel formulation; a Schaff-style Gaussian-mixture design search~\cite{schaff2019jointly} with a morphology-conditioned PPO-Lagrangian controller; one Lagrange multiplier per mixture component; and feasibility-first design ranking in the outer loop.
  \item \emph{Evidence that morphology is a safety lever.} On a hazard-field locomotion task the co-designed ant crosses about 55\% faster than the nominal ant at comparable constraint cost, and the evolved body changes with the hazard statistics. \HAVE (1 seed) \RUN seeds and controls.
  \item \emph{A failure mode of naive safe co-design.} With speed-only design ranking, the search returns a body 5$\times$ over budget when the constraint is hard to meet (random hazard fields). \HAVE \RUN the fix.
  \item \emph{An open-source GPU stack} for safe RL and safe co-design on MuJoCo Warp, with a different body in every environment lane. It runs 3.1--3.3$\times$ faster than the MJX-JAX implementation it extends on the same task (about 19$\times$ at equal solver accuracy), with physics equivalence checked; a 2B-step co-design run takes about 4\,h on one RTX 3090. \HAVE
  \end{enumerate}
\end{itemize}

\begin{figure}[t]
\centering
\figbox{4.2cm}{Teaser: path-traced frames of the nominal ant (top) and the co-designed ant (bottom) crossing the same hazard field at equal times.\\[3pt]
\HAVE renders: \texttt{renders/comparison/evolved\_vs\_nominal.mp4} and stills;\\ crossing time 9.3\,s (nominal) vs 5.4\,s (evolved) in these episodes.}
\caption{Same task, same safety budget, different body. \WRITE{} caption with the main numbers once final.}
\label{fig:teaser}
\end{figure}

%% =========================================================================
\section{Related Work}
\label{sec:related}
\sectarget{0.75 page}
\begin{itemize}
\item \textbf{Safe RL and its benchmarks.}
  \begin{itemize}
  \item CMDPs~\cite{altman1999constrained}; Lagrangian and primal-dual methods, including PPO-Lagrangian~\cite{ray2019benchmarking}.
  \item PID Lagrangian~\cite{stooke2020responsive}: relevant because we observed integral windup of the multiplier (Appendix~\ref{app:history}).
  \item CPO~\cite{achiam2017cpo}, CRPO~\cite{xu2021crpo}, Saut\'e RL~\cite{sootla2022saute}; surveys~\cite{garcia2015comprehensive,brunke2022safe}.
  \item Benchmarks: Safety Gym~\cite{ray2019benchmarking}, Safety-Gymnasium~\cite{ji2023safetygymnasium}, OmniSafe~\cite{ji2024omnisafe}.
  \item On accelerators: mjx-safety-gym~\cite{vanderhart_mjxsafetygym}, the MJX port our code extends, and CRAX~\cite{tomilin2026crax}: 8 tasks at 3 difficulty levels (Point, Ant, Humanoid and other agents), 100--200$\times$ faster than CPU benchmarks. Its physics steps on MJX-JAX; it already ships MuJoCo Warp for its pixel renderer (checked October 2026).
  \item All of these fix the body.
  \end{itemize}
\item \textbf{Constrained legged locomotion.} Constraints instead of reward shaping~\cite{kim2024constraints}, constraints as terminations~\cite{chanesane2024cat}, constrained visual parkour~\cite{chanesane2024soloparkour}, collision-free agile locomotion~\cite{he2024agile}. All on fixed bodies.
\item \textbf{Morphology--control co-design.}
  \begin{itemize}
  \item Evolutionary roots: virtual creatures~\cite{sims1994evolving}, physical robots~\cite{lipson2000automatic}, soft robots~\cite{cheney2013unshackling,bhatia2021evolution}.
  \item RL-based: learned design distributions~\cite{ha2019reinforcement,schaff2019jointly,schaff2022nlimb}, graph and transform policies~\cite{wang2019nge,yuan2022transform2act}, Darwinian evolution with learning~\cite{gupta2021embodied}, co-adaptation~\cite{luck2019data}, universal controllers~\cite{huang2020one,kurin2021my,gupta2022metamorph}.
  \item Constraints so far act on the \emph{design} (fixed parts and allowed ranges~\cite{constrainedevolution2022}; performance specifications~\cite{glamdring2026}) or enter as reward penalties~\cite{codesign2025notfree}; fall recovery as the objective~\cite{yue2025humanoid}.
  \item Literature search (2026-09-29, Semantic Scholar, arXiv and OpenAlex): of 911 papers citing 11 co-design seed papers, 57 mention safety and none trains the controller with constrained RL; of 2,552 papers citing 5 safe-RL seed papers, 21 mention design or morphology and none optimises a body. \CHECK{} update.
  \item Positioning: the first to put a budgeted, constrained-RL inner loop \emph{and} constraint handling into the design loop.
  \end{itemize}
\item \textbf{Constraint handling in evolutionary computation.}
  \begin{itemize}
  \item Penalties, feasibility rules~\cite{deb2000efficient}, stochastic ranking~\cite{runarsson2000stochastic}, $\varepsilon$-constrained methods~\cite{takahama2006constrained}; survey~\cite{coello2002theoretical}.
  \item Our outer loop is a search-distribution method in the EDA / natural-evolution-strategy family~\cite{larranaga2002estimation,wierstra2014natural}, updated by score-function gradients~\cite{williams1992simple}. Feasibility-first ranking is Deb's rule applied to its utilities.
  \item What is new: constraint values are noisy, non-stationary estimates, because the controller co-adapts with the design.
  \end{itemize}
\item \textbf{GPU simulation.}
  \begin{itemize}
  \item Brax~\cite{freeman2021brax}, Isaac Gym~\cite{makoviychuk2021isaac}, MJX and MuJoCo Playground~\cite{zakka2025playground}, MuJoCo Warp~\cite{mujocowarp2025}, Genesis~\cite{genesis2024}; massively parallel locomotion RL~\cite{rudin2022learning}.
  \item GPU simulators batch one kinematic structure, which is why our design space is parametric rather than topological (Sec.~\ref{sec:limits}).
  \end{itemize}
\end{itemize}

%% =========================================================================
\section{Problem Formulation}
\label{sec:problem}
\sectarget{0.5 page}
\begin{itemize}
\item \textbf{Design space.}
  \begin{itemize}
  \item $\xi \in \Xi = [0.6, 1.4]^7$: scale factors on the Safety Gym ant~\cite{ray2019benchmarking}, namely hip, thigh and shin lengths; hip, thigh and shin radii; and torso radius.
  \item Shared by all four legs. Nominal $\xi = \mathbf{1}$ (42.3\,kg).
  \item The kinematic structure is the same for every $\xi$ (a parametric design space).
  \end{itemize}
\item \textbf{Actuation.}
  \begin{itemize}
  \item Actuator gear is rescaled with mass $\times$ reach so the torque margin stays constant: we search over shape at constant relative strength.
  \item Without this, gear-to-gravity ratios spanned 0.94--3.82 and some bodies could not hold their own legs up.
  \item Bodies outside 15--90\,kg are rejected.
  \item \CHECK{} keep or change this (Limitations).
  \end{itemize}
\item \textbf{Fixed-design CMDP.} For a given $\xi$, $\max_\pi J_R(\pi;\xi)$ s.t.\ $J_C(\pi;\xi) \le d$, where $J_C$ is the expected \emph{total} cost of an episode and $d$ the budget.
\item \textbf{Safe co-design} (bilevel):
\begin{equation}
\max_{\xi \in \Xi}\; v(\pi^*_\xi;\xi) \quad \text{s.t.} \quad J_C(\pi^*_\xi;\xi) \le d,
\qquad \pi^*_\xi \in \arg\max_{\pi:\,J_C(\pi;\xi)\le d} J_R(\pi;\xi),
\label{eq:bilevel}
\end{equation}
with $v$ the crossing speed.
\item \textbf{Amortisation.} One policy $\pi_\theta(a \mid s, \xi)$, conditioned on the design~\cite{schaff2019jointly}, trained on bodies drawn from a search distribution $p_\phi(\xi)$, solves the inner problem for the whole population at once.
\item \textbf{Why speed and not return.}
  \begin{itemize}
  \item Return telescopes to distance travelled, so it saturates as soon as bodies arrive. In an early run all 8 bodies scored 22.49--22.50 of a 22.50 ceiling while their arrival times differed by 1.85$\times$.
  \item Design score $v$ = net forward displacement per decision. A non-arrival is charged the full 625-decision horizon, so ``dash two metres and fall'' never scores like a crossing.
  \end{itemize}
\end{itemize}

%% =========================================================================
\section{Method}
\label{sec:method}
\sectarget{1.25 pages including Fig.~\ref{fig:method} and Alg.~1}

\begin{figure*}[t]
\centering
\figbox{4.0cm}{Method diagram \WRITE: GMM design distribution $p_\phi$ $\rightarrow$ 16 designs per iteration $\rightarrow$ 1024 GPU lanes (64 per body) $\rightarrow$ morphology-conditioned PPO-Lagrangian with one $\lambda_k$ per mixture component $\rightarrow$ per-design (speed, cost) $\rightarrow$ feasibility-first ranking $\rightarrow$ REINFORCE update of $p_\phi$; chopping of mixture components at 400M / 600M / 800M; final fine-tune on the mode.}
\caption{Constrained co-design. Inner loop: a constrained controller for a population of bodies. Outer loop: an evolving design distribution that ranks designs feasibility-first.}
\label{fig:method}
\end{figure*}

\subsection{Population training on one GPU}
\begin{itemize}
\item Each design iteration samples $K = 16$ designs from $p_\phi$ and spreads them over 1024 parallel environments in blocks of 64 lanes per design.
\item Per-lane MuJoCo models differ only in 14 batched fields (geometry, inertia, gear). The structure is identical, so one compiled step function serves every body: no recompilation between iterations. Batched and individually compiled physics agree to about $10^{-8}$.
\item The design genes are appended to the observation (243 + 7 = 250 inputs).
\end{itemize}

\subsection{Inner loop: per-component PPO-Lagrangian}
\begin{itemize}
\item PPO~\cite{schulman2017proximal} with a separate cost critic; Lagrangian relaxation as in PPO-Lagrangian~\cite{ray2019benchmarking}:
\begin{equation}
\min_{\lambda \ge 0}\max_\theta\; \mathbb{E}_{k,\;\xi\sim p_\phi^{k}}\Big[J_R(\pi_\theta;\xi) - \lambda_k\big(J_C(\pi_\theta;\xi) - d\big)\Big].
\end{equation}
\item \textbf{One multiplier per mixture component $k$} (contribution). Each sample's cost advantage is weighted by $\lambda_{k(\xi)}$, so the budget has to hold for each subpopulation, not only on average.
\item Update: $\lambda_k \leftarrow \operatorname{clip}\big(\lambda_k + \eta_\lambda(\hat J_C^{k} - d),\,0,\,\lambda_{\max}\big)$ once per minibatch, with $\eta_\lambda = 1.5\times10^{-5}$, $\lambda_{\max} = 3$ (anti-windup) and $\lambda_0 = 0.01$.
\item Why the cap: without it, the integral-only multiplier ratcheted up and collapsed a working gait in earlier runs.
\item The budget is normalised by the episode length actually observed, so $d$ is an episode total and a policy cannot gain budget by loitering.
\item Flips and corridor exits end the episode with a one-off cost of 50. This closes the ``end the episode to escape future cost'' loophole that we observed with CRPO.
\end{itemize}

\subsection{Outer loop: an evolving design distribution}
\begin{itemize}
\item \textbf{Search distribution (following Schaff et al.~\cite{schaff2019jointly}).}
  \begin{itemize}
  \item A mixture of 8 diagonal Gaussians on $[-1, 1]^7$: means drawn from $U(-0.8, 0.8)$, standard deviation 0.577.
  \item Mixture weights frozen and uniform.
  \end{itemize}
\item \textbf{Update.}
  \begin{itemize}
  \item Score-function gradient through the sampled component only, evaluated at the unclipped sample $\tilde\xi$:
  \begin{equation}
  \nabla_\phi \approx \tfrac{1}{K}\textstyle\sum_{i=1}^{K} u_i\,\nabla_\phi \log p_\phi^{k_i}(\tilde\xi_i),
  \end{equation}
  where $u_i$ are scores standardised across the iteration's designs (current) or rank utilities (proposed, Sec.~\ref{sec:ffr}).
  \item Adam with learning rate 0.03.
  \item Designs with no completed episode are masked out.
  \end{itemize}
\item \textbf{Commitment by chopping.} The worst half of the live components is retired at 400M, 600M and 800M steps (8$\to$4$\to$2$\to$1), after a 200M warm-up with the distribution frozen. The last 300M steps fine-tune on the mode.
\item \textbf{Read as an evolutionary algorithm:} an EDA / NES-like search distribution~\cite{larranaga2002estimation,wierstra2014natural} with truncation selection between subpopulations.
\item \textbf{Sizing rule.} Each design iteration must contain at least one episode per lane: we use 3.85M steps, about 1.5 episodes. With shorter iterations only flips end episodes early, so scoring is biased against stable bodies.
\end{itemize}

\subsection{Constraint-aware design ranking (proposed)}
\label{sec:ffr}
\begin{itemize}
\item \RUN{} implement and evaluate.
\item Score each design by $(\hat v_i, \hat c_i)$: the mean speed and mean episode cost of its lanes in the iteration (about 96 episodes per design).
\item \textbf{Feasibility first} (Deb's rules~\cite{deb2000efficient}). Design $i$ ranks above $j$ if:
  \begin{itemize}
  \item $\hat c_i \le d < \hat c_j$; or
  \item both are feasible and $\hat v_i > \hat v_j$; or
  \item both are infeasible and $\hat c_i < \hat c_j$.
  \end{itemize}
\item Replace the standardised speeds with centred rank utilities (fitness shaping as in NES~\cite{wierstra2014natural}). Use the same ordering when chopping.
\item \textbf{Guards.}
  \begin{enumerate}
  \item Minimum progress: a body that does not arrive cannot be ``feasible'', so immobility is never rewarded.
  \item Noise: a confidence bound $\hat c_i + z\,\mathrm{SE}_i \le d$, or stochastic ranking~\cite{runarsson2000stochastic}.
  \end{enumerate}
\item \textbf{Ablation set:} speed only (Schaff, current); static penalty $v - \mu \max(0, c - d)$; feasibility first. \OPT a learned per-design multiplier $\lambda(\xi)$.
\end{itemize}

\begin{figure}[t]
\small
\fbox{\begin{minipage}{0.95\linewidth}
\textbf{Algorithm 1: one design iteration} \WRITE{} typeset properly (e.g.\ algorithm2e)
\begin{enumerate}
\item Sample $K$ designs $\tilde\xi_i \sim p_\phi$ (component $k_i$); clip into $\Xi$; build per-lane models.
\item Reset all lanes (redrawing random hazard fields if enabled).
\item Run PPO-Lagrangian epochs on the population. Update $\theta$, the critics and $\lambda_{k}$.
\item From completed episodes, estimate $\hat v_i$ and $\hat c_i$ for each design.
\item Rank the designs (speed only, or feasibility first) and compute utilities $u_i$.
\item Take a REINFORCE step on $\phi$. If a chop is due, retire the worst half of the components and reset the design optimiser.
\end{enumerate}
\end{minipage}}
\end{figure}

\subsection{Implementation}
\begin{itemize}
\item JAX with brax-style PPO; MuJoCo Warp~\cite{mujocowarp2025} through the MJX interface.
\item Everything runs on the device except the design update, which runs on the host between epochs.
\item Safe-RL algorithm implementations (PPO-Lagrangian, CRPO, Saut\'e) adapted from the ss2r and safe-learning codebases. \CHECK{} citations.
\item Details and measurements in Sec.~\ref{sec:res-sim} and Appendix~\ref{app:equiv}.
\end{itemize}

%% =========================================================================
\section{Experimental Setup}
\label{sec:setup}
\sectarget{1 page including Tables~\ref{tab:conditions} and \ref{tab:hparams}}

\subsection{Task: Minefield}
\begin{itemize}
\item \textbf{Corridor.} 11\,m long (start $x = -5$, finish line at $x = +6$) and 2\,m wide. Leaving it ends the episode with an exit cost of 50.
\item \textbf{Hazards.}
  \begin{itemize}
  \item 20 discs of radius 0.16\,m in the 10\,m $\times$ 2\,m band.
  \item Staggered lattice: 10 columns at 1\,m pitch, two discs per column, alternate columns shifted by $\pm 0.21$\,m.
  \item Fixed spawn at the corridor start, so on the lattice every episode has the same layout.
  \item Hazards do not collide: they are cost regions, as in Safety Gym.
  \end{itemize}
\item \textbf{Variants.}
  \begin{itemize}
  \item Radius 0.23\,m: 2.07$\times$ the area, covering 16.6\% of the band instead of 8.0\%.
  \item Random placement: rejection-sampled at every reset, so a 2B-step run sees about 530k distinct fields. Here safety cannot be a memorised stride.
  \end{itemize}
\item \textbf{The constraint is feasible.} A kinematic check shows a zero-cost crossing exists for the \emph{nominal} ant on the lattice: 77\% of torso positions in the band admit a cost-free four-foot stance, and they connect start to finish. Residual cost therefore reflects what the policy learned, not infeasibility.
\end{itemize}

\subsection{Robot and observations}
\begin{itemize}
\item \textbf{Robot.} The Safety Gym ant: 8 actuators, 42.3\,kg, condim 3.
\item \textbf{Timing.} RK4 with $\Delta t = 0.01$\,s and 2 substeps; action repeat 4, so one decision every 0.08\,s. An episode is at most 625 decisions (50\,s).
\item \textbf{Observation (243 values).}
  \begin{itemize}
  \item Proprioception and goal sensing (31).
  \item A 16-bin hazard lidar.
  \item A per-foot 7$\times$7 grid giving the cost that foot would incur if placed at each cell ($4 \times 49$).
  \item \CHECK{} the breakdown of the 31.
  \end{itemize}
\item \textbf{Co-design input.} The 7 design genes are added for co-design.
\item \textbf{Why the foot grid is there:} without it the Lagrangian could not reduce cost by improving foot placement (Appendix~\ref{app:history}).
\end{itemize}

\subsection{Reward, cost and budget}
\begin{itemize}
\item \textbf{Reward.} Forward progress $\Delta x$, plus goal-distance shaping, plus an upright bonus of 0.0002 per step. The episode ends on arrival, flip, exit or after 625 decisions.
\item \textbf{Cost per step.}
  \begin{itemize}
  \item For each hazard: linear penetration $\operatorname{clip}(1 - \delta/R, 0, 1)$ of the closest \emph{grounded} body part, where $\delta$ is the surface distance to the hazard centre and ``grounded'' means the part's lowest point is within 2\,cm of the floor. Summed over hazards.
  \item Plus a one-off 50 for a flip or exit.
  \item Graded like the hazard costs of Safety Gym and CRAX~\cite{ray2019benchmarking,tomilin2026crax}.
  \end{itemize}
\item \textbf{Budget.} $d$ is the expected episode total. The main experiments use $d = 2$, a strict budget.
\item \CHECK{} report what a careless crossing costs on the 2\,cm scale, to put $d = 2$ in context.
\end{itemize}

\subsection{Training and evaluation}
\begin{itemize}
\item \textbf{Training.} Hyperparameters in Table~\ref{tab:hparams}. Co-design runs take 2B steps: 520 design iterations of 3.85M steps. Single-body runs take 250--300M steps.
\item \textbf{Evaluation.}
  \begin{itemize}
  \item 512 stochastic episodes (128 environments $\times$ 4) on held-out layouts.
  \item Metrics: return, speed (m/s), cost (per-environment mean $\pm$ SE, $n = 128$), hazard steps (steps with any grounded part inside a disc), flipped \%, out-of-bounds \%, and the share of environments under budget.
  \end{itemize}
\item \textbf{Pitfall to state.} The trainer's periodic evaluations use 16 fixed random bodies, not the evolved one. Final bodies are evaluated separately with \path{scripts/eval_design.py}.
\item \textbf{Seeds.} 3 per co-design condition and 5 per single-body condition. \RUN
\item \textbf{Hardware.} RTX 3090 (cluster) and an RTX 4050 laptop. A 2B-step co-design run takes about 4--5\,h on Warp, against 13.9\,h on MJX-JAX.
\end{itemize}

\begin{table*}[t]
\caption{Experimental conditions. Hazards: L = lattice, R = random placement; radius in metres. Unless noted, the outer loop ranks designs by speed and the inner loop is PPO-Lagrangian with one $\lambda$ per mixture component.}
\label{tab:conditions}
\footnotesize
\setlength{\tabcolsep}{4pt}
\begin{tabular}{lllllll}
\toprule
ID & Body & Inner loop & Outer ranking & Hazards & Question & Status \\
\midrule
N & nominal, fixed & PPO-Lag, $d \in \{0, 2, 5\}$ & -- & L 0.16 & fixed-body baseline & $d=0,5$ \HAVE $d=2$ \RUN \\
CD & co-designed & PPO-Lag, $d=2$ & speed & L 0.16 & main method & \HAVE 1 seed \\
CD-0.23 & co-designed & PPO-Lag, $d=2$ & speed & L 0.23 & hazard size & \HAVE 1 seed \\
CD-R & co-designed & PPO-Lag, $d=2$ & speed & R 0.16 & hard constraint & \HAVE 1 seed \\
CD-R-FF & co-designed & PPO-Lag, $d=2$ & feasibility first & R 0.16 & outer-loop constraints & \RUN \\
E & evolved, fixed & PPO-Lag, $d=2$ (scratch) & -- & L 0.16 & body vs.\ co-training & \RUN \\
U & co-designed & none (unconstrained) & speed & L 0.16 & does safety shape the body? & \RUN \\
U$\to$S & U's body, fixed & PPO-Lag, $d=2$ & -- & L 0.16 & evolve, then add safety & \RUN \\
CD-shared & co-designed & PPO-Lag, one $\lambda$ & speed & L 0.16 & per-component $\lambda$ & \RUN \\
Sweep & both & PPO-Lag, $d \in \{0, 5, 25\}$ & speed or FF & L 0.16 & frontier (headline) & \RUN \\
UNI & uniform designs & PPO-Lag, $d=2$ & none; pick best & L 0.16 & is design learning needed? & \OPT \\
\bottomrule
\end{tabular}
\end{table*}

\begin{table}[t]
\caption{Hyperparameters (all runs). \CHECK{} against the launchers before submission.}
\label{tab:hparams}
\footnotesize
\begin{tabular}{ll}
\toprule
Parallel environments & 1024 (64 per design, 16 designs) \\
Policy / value networks & MLP $256\times4$ / $256\times5$ \\
Learning rate, entropy cost, clip & $10^{-4}$, $10^{-4}$, 0.3 \\
Discount (reward / cost) & 0.97 / 0.9 \\
Unroll, batch, minibatches & 10, 32, 32 \\
Multiplier lr / cap / initial value & $1.5\times10^{-5}$ / 3.0 / 0.01 \\
Flip / exit cost & 50 / 50 \\
Design components / lr / std init & 8 / 0.03 (Adam) / 0.577 \\
Design iteration length & 3.85M steps ($\approx$1.5 episodes per lane) \\
Warm-up / chops / fine-tune & 200M / 400, 600, 800M / last 300M \\
Total steps (co-design / single body) & 2B / 250--300M \\
Matmul precision & float32 \\
\bottomrule
\end{tabular}
\end{table}

%% =========================================================================
\section{Results}
\label{sec:results}
\sectarget{2.5 pages}

\subsection{Morphology improves the safety--performance trade-off}
\label{sec:res-main}
\begin{itemize}
\item \textbf{Headline (Table~\ref{tab:main}).} At budget 2 the evolved body crosses at 1.89\,m/s with cost $2.55 \pm 0.34$, and 66\% of environments stay under budget. The nominal ant reaches 1.20\,m/s at cost 1.9 ($d=0$) and 1.22\,m/s at cost 2.7 ($d=5$). That is about 55\% faster at comparable total cost. \HAVE (pre-Warp, 1 seed)
\item \textbf{Teaser numbers.} In the rendered episodes the crossing takes 5.4\,s against 9.3\,s (Fig.~\ref{fig:teaser}).
\item \textbf{State the claim precisely.} Today it is ``faster at matched cost'', not ``safer''.
  \begin{itemize}
  \item The evolved body's \emph{hazard} cost alone is about 1.48, against 0.65 ($d=0$) and 1.1 ($d=5$) for the nominal ant.
  \item It spends less of its budget on flips and exits (1.07, against 1.25 and 1.6).
  \item ``Safer'' needs the frontier: lower cost at equal speed, from the budget sweep (Fig.~\ref{fig:frontier}).
  \end{itemize}
\item \textbf{Confounds to close.}
  \begin{enumerate}
  \item Seeds. \RUN
  \item The nominal ant at exactly $d = 2$. \RUN
  \item Training budget: the nominal ant had 250M single-body steps; co-design had 2B steps in total, ending with a 300M fine-tune on the evolved body. Condition E (evolved body trained from scratch) separates ``better body'' from ``more training''. \RUN
  \item Evaluation protocol: the nominal numbers are means of the trainer's last three evaluations, the evolved ones come from the 512-episode protocol. Re-evaluate the nominal checkpoints with the same script. \CHECK
  \end{enumerate}
\end{itemize}

\begin{table*}[t]
\caption{Evaluation at budget $d = 2$ on the lattice field (radius 0.16). Cost is the per-environment episode total ($\pm$ SE, $n=128$); ``under budget'' is the share of environments whose mean cost is $\le d$. \HAVE rows are from the pre-Warp stack with one seed and must be regenerated. \CHECK{} the nominal rows are means of the trainer's last three evaluations, not the 512-episode protocol, and the $d=5$ run had not converged.}
\label{tab:main}
\footnotesize
\setlength{\tabcolsep}{4pt}
\begin{tabular}{llcccccccl}
\toprule
Body & Training & Return & Speed (m/s) & Cost & Hazard steps & Flipped (\%) & OOB (\%) & Under budget (\%) & Status \\
\midrule
Nominal & PPO-Lag $d=0$, 250M & 11.01 & 1.20 & 1.9 & 3.5 & 1.4 & 1.1 & -- & \HAVE \CHECK \\
Nominal & PPO-Lag $d=5$, 250M & 10.99 & 1.22 & 2.7 & 5.1 & 2.1 & 1.1 & -- & \HAVE \CHECK \\
Nominal & PPO-Lag $d=2$ & & & & & & & & \RUN \\
Co-designed & CD, 2B & 11.09 & 1.89 & $2.55 \pm 0.34$ & 6.6 & 0.4 & 1.8 & 66 & \HAVE \\
Co-designed, fixed & E, from scratch & & & & & & & & \RUN \\
Unconstrained design & U$\to$S & & & & & & & & \RUN \\
\bottomrule
\end{tabular}
\end{table*}

\begin{figure}[t]
\centering
\figbox{5.0cm}{Headline figure \RUN: speed (m/s, $y$) against episode cost ($x$, log scale) for the nominal ant across budgets $d \in \{0, 2, 5, 25\}$ (the fixed-body frontier) and for co-designed bodies at the same budgets; budget lines dashed; error bars over seeds. The claim holds if the co-design frontier dominates.}
\caption{Safety--performance frontier: fixed body vs.\ co-designed body.}
\label{fig:frontier}
\end{figure}

\subsection{What evolves}
\label{sec:res-what}
\begin{itemize}
\item \textbf{Lattice, radius 0.16} (Table~\ref{tab:genes}). Hip, thigh and shin lengths $\times$1.40 and hip radius $\times$1.40, all at the box edge; thigh radius $\times$0.94, shin $\times$0.83, torso $\times$0.72. Mass 38.5\,kg (nominal 42.3); gear 191 (nominal 150). \HAVE
\item \textbf{Consistent across all three regimes:} hip and thigh at the length limit, thick hips, a small torso.
\item \textbf{What changes:} shin length, 1.40 (lattice 0.16) $\to$ 1.22 (lattice 0.23) $\to$ 1.03 (random). One seed each, so a hint only.
\item \textbf{Interpretation} (to test):
  \begin{itemize}
  \item Long legs give clearance to step over discs, and longer strides.
  \item A small torso lowers mass and inertia.
  \item On random fields, shorter shins may favour precise placement over clearance.
  \end{itemize}
\item \textbf{Gait analysis to back it} \RUN: share of time the torso is above a disc, foot-placement histograms relative to discs, step length.
\item \textbf{Weakness to own:} 4 of 7 genes end at the upper bound, so the bounds bind (Sec.~\ref{sec:limits}).
\item \textbf{Fig.~\ref{fig:genes}:} the mode of each component across training, with chop events. \HAVE design logs; \WRITE plot.
\end{itemize}

\begin{table}[t]
\caption{Evolved designs (scale factors on the nominal ant). Hip = \texttt{aux}, thigh = \texttt{leg}, shin = \texttt{ankle} segments of the Safety Gym ant. \HAVE (pre-Warp, 1 seed).}
\label{tab:genes}
\footnotesize
\setlength{\tabcolsep}{3pt}
\begin{tabular}{lccccccccc}
\toprule
 & \multicolumn{3}{c}{Length} & \multicolumn{3}{c}{Radius} & Torso & Mass & Gear \\
\cmidrule(lr){2-4}\cmidrule(lr){5-7}
Run & hip & thigh & shin & hip & thigh & shin & radius & (kg) & \\
\midrule
Nominal & 1.00 & 1.00 & 1.00 & 1.00 & 1.00 & 1.00 & 1.00 & 42.3 & 150 \\
L 0.16 & 1.40 & 1.40 & 1.40 & 1.40 & 0.94 & 0.83 & 0.72 & 38.5 & 191 \\
L 0.23 & 1.40 & 1.40 & 1.22 & 1.31 & 0.76 & 0.96 & 0.60 & 46.2 & 219 \\
R 0.16 & 1.40 & 1.40 & 1.03 & 1.32 & 0.82 & 0.99 & 0.80 & 42.5 & 193 \\
\bottomrule
\end{tabular}
\end{table}

\begin{figure}[t]
\centering
\figbox{4.0cm}{Gene trajectories \WRITE: mode of each mixture component across training (7 panels, or 2D projections), with chop events at 400/600/800M and the box limits marked. Data: the design lines in the co-design logs and the \texttt{design\_state.npz} checkpoints.}
\caption{How the design distribution moves and commits.}
\label{fig:genes}
\end{figure}

\subsection{The outer loop needs constraint handling}
\label{sec:res-outer}
\begin{itemize}
\item \textbf{Random hazard fields} (Table~\ref{tab:regimes}). The evolved body crosses at 1.48\,m/s but at cost $10.15 \pm 0.60$, 5$\times$ the budget. No environment stays under budget. \HAVE
\item \textbf{Training trace.} Every component's $\lambda$ was pinned at the cap (3.0), with the constraint at $-0.71$: the inner loop could not meet $d = 2$ for any subpopulation. \HAVE
\item \textbf{Mechanism.} Speed-only ranking kept favouring the fastest designs, which were also the most violating. The search had no reason to prefer a feasible body.
\item \textbf{On the lattice} the same loop happened to settle near the budget ($2.55 \pm 0.34$ and $2.31 \pm 0.20$; both $+1.5$--$1.6$ SE over $d$, not significant at 2 SE). Feasibility there came from the inner loop, not from the search.
\item \textbf{The fix:} feasibility-first ranking, CD-R vs CD-R-FF, 3 seeds each (Fig.~\ref{fig:outer}). Report the final body, the cost trace, and whether the search finds a feasible body or proves none is learnable. \RUN
\item \textbf{The per-component $\lambda$ ablation} (CD vs CD-shared) also belongs here: with one shared $\lambda$, does a reckless component get carried? \RUN
\end{itemize}

\begin{table*}[t]
\caption{Evolved bodies under three hazard regimes, budget $d = 2$, 512 episodes each. \HAVE rows: pre-Warp stack, one seed.}
\label{tab:regimes}
\footnotesize
\begin{tabular}{llcccccccl}
\toprule
Hazards & Outer ranking & Return & Speed (m/s) & Cost & Hazard steps & Flipped (\%) & OOB (\%) & Under budget (\%) & Status \\
\midrule
Lattice, $r=0.16$ & speed & 11.09 & 1.89 & $2.55 \pm 0.34$ & 6.6 & 0.4 & 1.8 & 66 & \HAVE \\
Lattice, $r=0.23$ & speed & 11.21 & 1.63 & $2.31 \pm 0.20$ & 25.3 & 0.4 & 0.0 & 60 & \HAVE \\
Random, $r=0.16$ & speed & 10.88 & 1.48 & $10.15 \pm 0.60$ & 25.8 & 1.4 & 4.9 & 0 & \HAVE \\
Random, $r=0.16$ & feasibility first & & & & & & & & \RUN \\
\bottomrule
\end{tabular}
\end{table*}

\begin{figure}[t]
\centering
\figbox{4.0cm}{\RUN: episode cost of the designs sampled at each iteration (median and IQR) against training steps, speed-only vs feasibility-first on random fields, with the budget line; inset: per-component $\lambda$ (pinned at the cap under speed-only).}
\caption{Speed-only design ranking drifts to infeasible bodies; feasibility-first ranking does not.}
\label{fig:outer}
\end{figure}

\subsection{Does the safety constraint shape the body?}
\label{sec:res-shape}
\begin{itemize}
\item This is the core of the ``lever'' claim. If unconstrained co-design evolves the same body, safety is a by-product of speed, not a design driver.
\item \textbf{Compare} U against CD (genes and performance), U$\to$S against E (is evolving \emph{under} the constraint better than evolving first and adding safety later?), and the evolved genes against the budget across the sweep. \RUN
\item \textbf{Honest expectation.} On this task speed and safety probably pull the same way: long legs help both. The bodies may differ little, which motivates a task with opposing pressures (Sec.~\ref{sec:discussion}).
\end{itemize}

\subsection{Simulation stack: throughput and fidelity}
\label{sec:res-sim}
\begin{itemize}
\item \textbf{Throughput} (Table~\ref{tab:throughput}). On an RTX 3090 at 1024 environments:
  \begin{itemize}
  \item 38.2k $\to$ 118.3k steps/s for a single body (3.1$\times$) and 40.4k $\to$ 133.3k for co-design with 16 bodies (3.3$\times$), against the MJX-JAX backend with its solver capped for speed. \HAVE
  \item At equal solver accuracy (both run to convergence) the laptop shows 3.8k $\to$ 73.3k training steps/s (19$\times$), and 4.4k $\to$ 140k steps/s for pure stepping (32$\times$). \HAVE \CHECK{} repeat on the 3090
  \end{itemize}
\item \textbf{Why it is faster.} MJX-JAX pays for every \emph{possible} contact because of static array shapes (about 188 constraint rows per environment here), and its line search runs a fixed number of steps. Warp only processes active contacts (at most 20 rows per environment) and stops once every world has converged~\cite{zakka2025playground,mujocowarp2025}.
\item \textbf{Scaling} (Fig.~\ref{fig:scaling}): 4096 environments give 220k (single) and 233k (co-design) steps/s, 1.75$\times$ the 1024-environment rate. Learning at that size is unverified, and co-design iterations must lengthen to keep at least one episode per lane. \HAVE
\item \textbf{Fidelity} (Appendix~\ref{app:equiv}).
  \begin{itemize}
  \item Old and new MuJoCo agree to float32 rounding once the solver converges.
  \item Warp and JAX agree to about $10^{-7}$, including with a different body in every lane.
  \item The vase cost, re-implemented with contact sensors, matches the original step for step.
  \item A 300M-step learning check reproduces the old baseline: 1.52 vs 1.45\,m/s, cost 26.3 vs 27.4 at $d = 25$. \HAVE
  \end{itemize}
\item \textbf{Wall clock.} A 2B-step co-design run drops from 13.9\,h to about 4\,h of training, which is what makes seeds and sweeps affordable.
\item \textbf{Positioning} (decision pending).
  \begin{itemize}
  \item CRAX~\cite{tomilin2026crax} already provides a broad MJX safe-RL suite, including Ant tasks.
  \item The Warp backend alone is a temporary difference: CRAX already depends on MuJoCo Warp, and our own switch took about 300 changed lines in 10 files.
  \item Our claim is therefore neither breadth nor the backend. It is (i) a different body in every lane of one batch, with constraint costs computed on the GPU, (ii) the safe co-design loop, and (iii) validated equivalence of the Warp port, including per-lane bodies.
  \item \OPT{} same-task throughput: CRAX's Ant goal task against ours; Safety-Gymnasium on CPU.
  \end{itemize}
\end{itemize}

\begin{table}[t]
\caption{Training throughput (environment steps per second, compilation excluded). Single = one nominal body; co-design = 16 bodies. Solver: \emph{capped} = 4 iterations / 8 line-search steps (the MJX-JAX speed setting), \emph{full} = MuJoCo defaults, run to convergence. \CHECK{} re-measure all cells with one script on the 3090.}
\label{tab:throughput}
\footnotesize
\setlength{\tabcolsep}{3pt}
\begin{tabular}{lcccc}
\toprule
 & \multicolumn{2}{c}{RTX 4050 laptop} & \multicolumn{2}{c}{RTX 3090} \\
\cmidrule(lr){2-3}\cmidrule(lr){4-5}
Backend, solver & single & co-design & single & co-design \\
\midrule
MJX-JAX, full & 3.8k & -- & \RUN & -- \\
MJX-JAX, capped & 13.8k & 19.0k & 38.2k & 40.4k \\
Warp, full & 73.3k & 100.0k & 118.3k & 133.3k \\
Warp, full, 4096 envs & -- & -- & 220k & 233k \\
\bottomrule
\end{tabular}
\\[2pt]{\scriptsize Laptop single-body runs used 512 environments; all other cells 1024 unless noted.}
\end{table}

\begin{figure}[t]
\centering
\figbox{3.6cm}{\HAVE data, \WRITE plot: steps/s vs number of environments (1024--8192) on the 3090, Warp, single and co-design; MJX-JAX points for reference. GPU utilisation read 100\% at every size, so small batches under-fill the GPU.}
\caption{Throughput scaling with environments per GPU.}
\label{fig:scaling}
\end{figure}

%% =========================================================================
\section{Discussion}
\label{sec:discussion}
\sectarget{0.75 page}
\begin{itemize}
\item \textbf{Why the body helps.}
  \begin{itemize}
  \item The nominal ant \emph{can} cross at zero cost (Sec.~\ref{sec:setup}), but constrained RL does not find that precise foot placement.
  \item A body with more clearance and a longer stride widens the margin, so the safe crossing becomes easier to learn.
  \item Hypothesis: morphology helps safety mainly through \emph{learnability}, not feasibility. Test: rerun the kinematic stance check for the evolved bodies (cheap, CPU only). \RUN
  \end{itemize}
\item \textbf{Constraint handling is part of the method, not a detail.} The inner loop protects safety only when it succeeds; the outer loop has to know when it did not.
\item \textbf{Limitations.}
\label{sec:limits}
  \begin{enumerate}
  \item One seed per condition and a pre-Warp stack so far.
  \item Bounds bind: 4 of 7 genes end at the box edge.
  \item Size is free: actuators are rescaled, so bigger limbs cost no relative strength.
  \item Speed and safety point the same way on this task, so there is no tension to resolve.
  \item On the lattice, safety can be a memorised stride; random fields are the hard case, and they fail without feasibility-first ranking.
  \item Parametric design only. GPU simulators batch one kinematic structure, so topology needs a superset model or per-world kinematic trees (a draft MuJoCo Warp feature).
  \item Simulation only; no sim-to-real.
  \item Designs are scored through a shared policy, which generalises by interpolation (novel bodies inside the trained mass range: 7/7 worked) but not by extrapolation (outside it: 3/7).
  \end{enumerate}
\item \textbf{Future work.}
  \begin{itemize}
  \item A lasers + mines task: beams at random heights. Long legs step over mines but hit beams; low bodies duck beams but must weave between mines, so the best body depends on the mix and on the budget.
  \item A payload (tipping or jerk cost) and an energy cost.
  \item Topology through a superset ant.
  \item A real quadruped model (e.g.\ from MuJoCo Menagerie).
  \item Quality diversity: MAP-Elites over morphology with safety descriptors~\cite{mouret2015illuminating}; CMA-ES as an alternative outer loop~\cite{hansen2001completely}.
  \end{itemize}
\end{itemize}

%% =========================================================================
\section{Conclusion}
\label{sec:conclusion}
\sectarget{0.25 page}
\begin{itemize}
\item Safety is a property of the body as well as of the controller.
\item Co-designing morphology under a safety constraint found an ant that is about 55\% faster at comparable cost. \CHECK{} final wording once the frontier is in.
\item Safe co-design needs constraint handling at both levels: per-component multipliers inside, feasibility-first ranking outside.
\item A GPU stack makes this affordable: billions of steps in hours on one consumer GPU.
\end{itemize}

\begin{acks}
\WRITE{} Cluster and compute acknowledgements; funding. Removed automatically in the anonymous version.
\end{acks}

\bibliographystyle{ACM-Reference-Format}
\bibliography{references}

%% =========================================================================
\appendix
\noindent\textit{For the October report these sections are part of the main text. For GECCO they move to supplementary material.}

\section{Hyperparameters and reproducibility}
\label{app:repro}
\begin{itemize}
\item Full flag listing per condition, from the launchers (\path{cluster/ant_morphology_codesign.sbatch}: knobs BUDGET, HAZARD\_SIZE, HAZARD\_PLACEMENT, PENALIZER). \WRITE
\item \textbf{Code.} Public repository (anonymised for review). Tag \texttt{pre-warp} reproduces every \HAVE number; \texttt{main} is the Warp stack.
\item \textbf{Evaluation.} \path{scripts/eval_design.py} (evolved bodies), \path{scripts/eval_checkpoint.py} (single bodies).
\item \textbf{Checkpoints.} Run names: \path{ant_codesign_lagrangian_b2_grid7_2000M} and its \texttt{\_hz0.23} and \texttt{\_random} variants; \path{ant_minefield_final_eps2_b{0,5}_250M}.
\end{itemize}

\section{Development history and negative results}
\label{app:history}
\noindent\textit{Report only. The decisions that shaped the final recipe, each backed by a measurement.}
\begin{itemize}
\item \textbf{Goal-reaching was unlearnable for the ant.} With the goal in a random direction, a scripted gait that walks several metres earns $+0.078 \pm 0.672$ return over 16 seeds. Hence the corridor task, with reward linear in forward progress.
\item \textbf{Policy capacity mattered most early on.} Widening the policy from $32\times4$ to $256\times4$ cut cost 2.7$\times$ at the same reward.
\item \textbf{Constraint mechanisms first ``succeeded'' by doing less of the task.}
  \begin{itemize}
  \item CRPO: cost per step \emph{doubled} (0.025 $\to$ 0.050) while episode cost fell, because episodes got shorter (2035 $\to$ 738 steps).
  \item Saut\'e: the budget was not a lever (tightening it tenfold, 500 $\to$ 50, did not lower cost: 84.9 vs 95.7); penalty 5 cut cost 43\% and reward 42\% together.
  \item Lagrangian without a cap: the multiplier wound up and the policy met the budget by walking 3\,m and flipping.
  \end{itemize}
\item \textbf{Fixes, each measured.}
  \begin{itemize}
  \item A flip/exit cost.
  \item A graded penetration cost (Safety Gym / CRAX style).
  \item A budget normalised by the observed episode length.
  \item A multiplier cap.
  \item \emph{The per-foot hazard grid} (50M-step arms, one seed each). Without the grid, the Lagrangian ended \emph{above} its unconstrained control (cost 78.6 vs 63.1). With it, the Lagrangian cut cost below its unconstrained control: by 12\% from scratch (44.8 vs 51.1), and by 35\% when warm-started from that control (33.0 vs 51.1), with the crossing intact.
  \item A $7\times7$ grid beat $5\times5$: cost 43.3 vs 53.8, speed 1.08 vs 0.99.
  \end{itemize}
\item \textbf{Design scoring.} Return cannot rank bodies here (it saturates), hence the speed score. Iterations shorter than one episode per lane bias scoring towards bodies that flip early.
\item \textbf{Generalisation of the shared policy.} The boundary is mass, not distance in gene space: novel bodies inside the trained mass range worked (7/7, mean arrival 99\%), outside it 3/7 (mean 63\%).
\end{itemize}

\section{Throughput history: from the original mjx-safety-gym to now}
\label{app:speed}
\noindent\textit{Report only.} Our code started from mjx-safety-gym 0.1.6~\cite{vanderhart_mjxsafetygym}. Steps of the chain were measured on different machines, environment counts and task versions, so the factors do not multiply exactly. Only the end-to-end comparison is a direct before/after measurement.
\begin{table*}[t]
\caption{Throughput history of the ant pipeline (report only).}
\label{tab:history}
\footnotesize
\begin{tabular}{lll}
\toprule
When (2026) & Change & Measured effect \\
\midrule
25 Jul & JAX 0.6.0 $\to$ 0.10.2 & $\sim$2.4$\times$ MJX stepping \CHECK \\
9 Aug & first ant baseline (Goal task, vases, condim 6) & 668--798 steps/s (3090, 256 envs) \\
13 Aug & contact dimension 6 $\to$ 3 & 1.82$\times$ (laptop, 128 envs) \\
15 Aug & Minefield task (no vases; hazards non-colliding) & 3.13$\times$ (laptop, 256 envs) \\
17 Aug & 1024 environments (3090 knee) & +31\% over 512 \\
22 Aug--26 Sep & corridor walls added, then removed & $-27\%$, then +37\% \\
27 Sep & solver capped at 4/8 (MJX-JAX) & 3.5$\times$ (laptop) \\
4 Oct & MuJoCo Warp, full solver & 3.1--3.3$\times$ vs capped JAX (3090) \\
4 Oct & 4096 environments on Warp & 1.75$\times$ vs 1024 (3090) \\
\midrule
End to end & 3090, ant & $\sim$0.8k $\to$ 118--133k ($\approx$150$\times$); 220--233k at 4096 envs \\
\bottomrule
\end{tabular}
\end{table*}
\begin{itemize}
\item \textbf{Engineering, same physics:} the JAX upgrade, environment counts, the Warp backend.
\item \textbf{Changes to what is simulated:} condim 3 (Gym's setting), the minefield task (vases removed), walls, and the solver cap (an approximation, since replaced by Warp's converged solver).
\end{itemize}

\section{Simulator equivalence tests}
\label{app:equiv}
\begin{itemize}
\item \textbf{One step from identical states.}
  \begin{itemize}
  \item MuJoCo 3.3.2 vs 3.14 agree to float32 rounding when the solver converges.
  \item With the 4/8 cap they differ by up to $1.4\times10^{-3}$ in joint velocity.
  \item Warp vs JAX agree to about $10^{-7}$, including with per-lane bodies.
  \end{itemize}
\item \textbf{Vase cost via contact sensors} (MJX-JAX lacks subtree matching, so one sensor per vase--geom pair). It needs found \emph{and} negative distance: 0 differing steps against the contact-list cost, while ``found'' alone over-counts by about 9\% (5,735 vs 5,283).
\item \textbf{Policy transfer.} A policy trained under the capped solver is 4.5\% slower on the new stack (1.141 vs 1.195\,m/s, 12 SE); cost, flips and exits are unchanged. Hence: never mix stacks in one comparison.
\item \textbf{Learning check} (single nominal body, $d = 25$, 5\,cm gate, mean of evaluations at 241--261M steps):
\end{itemize}
\begin{center}
\footnotesize
\setlength{\tabcolsep}{3pt}
\begin{tabular}{lcccccc}
\toprule
Stack & Return & Speed & Cost & Hz.\ steps & Flip \% & OOB \% \\
\midrule
JAX capped (3090, 1.7\,h) & 10.89 & 1.45 & 27.4 & 70.8 & 2.6 & 1.9 \\
Warp full (laptop, 46\,min) & 11.06 & 1.52 & 26.3 & 76.6 & 0.3 & 2.4 \\
\bottomrule
\end{tabular}
\end{center}

\end{document}
